\documentclass[10pt,a4paper]{amsart}
\usepackage{amsmath,amssymb,amsthm,graphicx}
\usepackage{verbatim}
\newenvironment{reference}{\comment}{\endcomment}
%\newenvironment{reference}{}{}
\newenvironment{slogan}{\comment}{\endcomment}
\newenvironment{history}{\comment}{\endcomment}

% For commutative diagrams we use Xy-pic
\usepackage[all]{xy}

% We use 2cell for 2-commutative diagrams.
\xyoption{2cell}
% Theorem environments.
%
\theoremstyle{plain}
\newtheorem{theorem}[subsection]{Theorem}
\newtheorem{proposition}[subsection]{Proposition}
\newtheorem{lemma}[subsection]{Lemma}

\theoremstyle{definition}
\newtheorem{definition}[subsection]{Definition}
\newtheorem{example}[subsection]{Example}
\newtheorem{exercise}[subsection]{Exercise}
\newtheorem{situation}[subsection]{Situation}

\theoremstyle{remark}
\newtheorem{remark}[subsection]{Remark}
\newtheorem{remarks}[subsection]{Remarks}

\numberwithin{equation}{subsection}

% Macros
%
\def\lim{\mathop{\mathrm{lim}}\nolimits}
\def\colim{\mathop{\mathrm{colim}}\nolimits}
\def\Spec{\mathop{\mathrm{Spec}}}
\def\Hom{\mathop{\mathrm{Hom}}\nolimits}
\def\Ext{\mathop{\mathrm{Ext}}\nolimits}
\def\SheafHom{\mathop{\mathcal{H}\!\mathit{om}}\nolimits}
\def\SheafExt{\mathop{\mathcal{E}\!\mathit{xt}}\nolimits}
\def\Sch{\mathit{Sch}}
\def\Mor{\mathop{\mathrm{Mor}}\nolimits}
\def\Ob{\mathop{\mathrm{Ob}}\nolimits}
\def\Sh{\mathop{\mathit{Sh}}\nolimits}
\def\NL{\mathop{N\!L}\nolimits}
\def\CH{\mathop{\mathrm{CH}}\nolimits}
\def\proetale{{pro\text{-}\acute{e}tale}}
\def\etale{{\acute{e}tale}}
\def\QCoh{\mathit{QCoh}}
\def\Ker{\mathop{\mathrm{Ker}}}
\def\Im{\mathop{\mathrm{Im}}}
\def\Coker{\mathop{\mathrm{Coker}}}
\def\Coim{\mathop{\mathrm{Coim}}}

% Boxtimes
%
\DeclareMathSymbol{\boxtimes}{\mathbin}{AMSa}{"02}

%
% Macros for moduli stacks/spaces
%
\def\QCohstack{\mathcal{QC}\!\mathit{oh}}
\def\Cohstack{\mathcal{C}\!\mathit{oh}}
\def\Spacesstack{\mathcal{S}\!\mathit{paces}}
\def\Quotfunctor{\mathrm{Quot}}
\def\Hilbfunctor{\mathrm{Hilb}}
\def\Curvesstack{\mathcal{C}\!\mathit{urves}}
\def\Polarizedstack{\mathcal{P}\!\mathit{olarized}}
\def\Complexesstack{\mathcal{C}\!\mathit{omplexes}}
% \Pic is the operator that assigns to X its picard group, usage \Pic(X)
% \Picardstack_{X/B} denotes the Picard stack of X over B
% \Picardfunctor_{X/B} denotes the Picard functor of X over B
\def\Pic{\mathop{\mathrm{Pic}}\nolimits}
\def\Picardstack{\mathcal{P}\!\mathit{ic}}
\def\Picardfunctor{\mathrm{Pic}}
\def\Deformationcategory{\mathcal{D}\!\mathit{ef}}
\usepackage{cleveref}
\crefname{originalreference}{original reference}{original references}
\Crefname{originalreference}{Original reference}{Original references}
\makeatletter
\newlabel{homology-lemma-characterize-weak-serre-subcategory}{{0754}{source}{}{sourcepointer-1}{}}
\newlabel{homology-lemma-characterize-weak-serre-subcategory@cref}{{[originalreference][0][]0754}{[1][0][]source}{}{}{}}
\makeatother
\crefname{originallemma}{original lemma}{original lemmas}
\Crefname{originallemma}{Original lemma}{Original lemmas}
\makeatletter
\newlabel{lemma-lisse-etale-modules}{{0789}{source}{}{sourcepointer-2}{}}
\newlabel{lemma-lisse-etale-modules@cref}{{[originallemma][0][]0789}{[1][0][]source}{}{}{}}
\makeatother
\crefname{originallemma}{original lemma}{original lemmas}
\Crefname{originallemma}{Original lemma}{Original lemmas}
\makeatletter
\newlabel{lemma-shriek-quasi-coherent}{{07B0}{source}{}{sourcepointer-3}{}}
\newlabel{lemma-shriek-quasi-coherent@cref}{{[originallemma][0][]07B0}{[1][0][]source}{}{}{}}
\makeatother
\crefname{originalreference}{original reference}{original references}
\Crefname{originalreference}{Original reference}{Original references}
\makeatletter
\newlabel{stacks-sheaves-lemma-quasi-coherent-algebraic-stack}{{06WV}{source}{}{sourcepointer-4}{}}
\newlabel{stacks-sheaves-lemma-quasi-coherent-algebraic-stack@cref}{{[originalreference][0][]06WV}{[1][0][]source}{}{}{}}
\makeatother
\crefname{originallemma}{original lemma}{original lemmas}
\Crefname{originallemma}{Original lemma}{Original lemmas}
\makeatletter
\newlabel{lemma-lisse-etale}{{0788}{source}{}{sourcepointer-5}{}}
\newlabel{lemma-lisse-etale@cref}{{[originallemma][0][]0788}{[1][0][]source}{}{}{}}
\makeatother
\crefname{originallemma}{original lemma}{original lemmas}
\Crefname{originallemma}{Original lemma}{Original lemmas}
\makeatletter
\newlabel{lemma-quasi-coherent}{{07B1}{source}{}{sourcepointer-6}{}}
\newlabel{lemma-quasi-coherent@cref}{{[originallemma][0][]07B1}{[1][0][]source}{}{}{}}
\makeatother
\crefname{originalproposition}{original proposition}{original propositions}
\Crefname{originalproposition}{Original proposition}{Original propositions}
\makeatletter
\newlabel{proposition-loc-qcoh-flat-base-change}{{0771}{source}{}{sourcepointer-7}{}}
\newlabel{proposition-loc-qcoh-flat-base-change@cref}{{[originalproposition][0][]0771}{[1][0][]source}{}{}{}}
\makeatother
\crefname{originallemma}{original lemma}{original lemmas}
\Crefname{originallemma}{Original lemma}{Original lemmas}
\makeatletter
\newlabel{lemma-lisse-etale-cohomology}{{0GR0}{source}{}{sourcepointer-8}{}}
\newlabel{lemma-lisse-etale-cohomology@cref}{{[originallemma][0][]0GR0}{[1][0][]source}{}{}{}}
\makeatother
\crefname{originalreference}{original reference}{original references}
\Crefname{originalreference}{Original reference}{Original references}
\makeatletter
\newlabel{descent-proposition-equivalence-quasi-coherent}{{03DX}{source}{}{sourcepointer-9}{}}
\newlabel{descent-proposition-equivalence-quasi-coherent@cref}{{[originalreference][0][]03DX}{[1][0][]source}{}{}{}}
\makeatother
\crefname{originalreference}{original reference}{original references}
\Crefname{originalreference}{Original reference}{Original references}
\makeatletter
\newlabel{descent-proposition-same-cohomology-quasi-coherent}{{03DW}{source}{}{sourcepointer-10}{}}
\newlabel{descent-proposition-same-cohomology-quasi-coherent@cref}{{[originalreference][0][]03DW}{[1][0][]source}{}{}{}}
\makeatother
\crefname{originalreference}{original reference}{original references}
\Crefname{originalreference}{Original reference}{Original references}
\makeatletter
\newlabel{coherent-lemma-quasi-coherent-affine-cohomology-zero}{{01XB}{source}{}{sourcepointer-11}{}}
\newlabel{coherent-lemma-quasi-coherent-affine-cohomology-zero@cref}{{[originalreference][0][]01XB}{[1][0][]source}{}{}{}}
\makeatother
\title{Quasi-coherent modules are a weak Serre subcategory on the lisse-\textbackslash{}'etale and flat-fppf sites}
\author{The Stacks Project Authors}
\date{stacks-project:stacks-cohomology.tex:local-snapshot-20261012; GNU FDL 1.2 or later}
\begin{document}
\maketitle
\noindent\textit{Editorial note.} All delivered mathematical text is the author's own, copied verbatim from the pinned Stacks Project chapter file stacks-cohomology.tex; no statement or proof step was rewritten, reordered or removed.\par\medskip
\noindent\textit{Editorial note.} The counted span cites items from Homology, Sheaves on Stacks, Descent and Cohomology of Schemes plus same-chapter lemmas, recorded in assembly.external\_references with canonical tags as navigation pointers.\par\medskip
\noindent\textit{Allowed standard prerequisites (editorial).}\begin{itemize}
\item Weak Serre subcategory criterion. Stacks Project, Homology, Lemma homology-lemma-characterize-weak-serre-subcategory. Supplies the four conditions verified.
\item \$g\_!\$-\$g\textasciicircum{}\$* transport on the lisse-\textbackslash{}'etale site. Stacks Project, Cohomology of Algebraic Stacks, Lemmas lemma-lisse-etale-modules, lemma-shriek-quasi-coherent, lemma-lisse-etale, lemma-quasi-coherent. Supplies \$g\textasciicircum{}* g\_! = \textbackslash{}mathrm\{id\}\$, quasi-coherence of \$g\_!\$, exactness of \$g\textasciicircum{}*\$, and pullback identification.
\item Cokernels of quasi-coherent maps are quasi-coherent. Stacks Project, Sheaves on Stacks, Lemma stacks-sheaves-lemma-quasi-coherent-algebraic-stack. Gives the cokernel half of condition (3).
\item Kernels in LQCoh\textasciicircum{}\{fbc\}. Stacks Project, Cohomology of Algebraic Stacks, Proposition proposition-loc-qcoh-flat-base-change. Gives the kernel half of condition (3).
\item Vanishing of \$H\textasciicircum{}1\$ of quasi-coherent modules on affine schemes. Stacks Project, Descent, Propositions descent-proposition-equivalence-quasi-coherent and descent-proposition-same-cohomology-quasi-coherent; Cohomology of Schemes, Lemma coherent-lemma-quasi-coherent-affine-cohomology-zero. Assembles the \$H\textasciicircum{}1 = 0\$ used to lift the presentation.
\end{itemize}\medskip
\ifdefined\hypertarget\hypertarget{sourcepointer-1}{}\fi
\noindent\textit{Source reference.} Original reference 0754: sources/stacks-chow/tags.txt, source lines 9118--9118. Cross-chapter standard prerequisite cited by the author inside the counted span; it is used as a pointer and its own proof is not part of this item.\par\medskip
\ifdefined\hypertarget\hypertarget{sourcepointer-2}{}\fi
\noindent\textit{Source reference.} Original lemma 0789: sources/stacks-chow/tags.txt, source lines 9229--9229. Cross-chapter standard prerequisite cited by the author inside the counted span; it is used as a pointer and its own proof is not part of this item.\par\medskip
\ifdefined\hypertarget\hypertarget{sourcepointer-3}{}\fi
\noindent\textit{Source reference.} Original lemma 07B0: sources/stacks-chow/tags.txt, source lines 9325--9325. Cross-chapter standard prerequisite cited by the author inside the counted span; it is used as a pointer and its own proof is not part of this item.\par\medskip
\ifdefined\hypertarget\hypertarget{sourcepointer-4}{}\fi
\noindent\textit{Source reference.} Original reference 06WV: sources/stacks-chow/tags.txt, source lines 8829--8829. Cross-chapter standard prerequisite cited by the author inside the counted span; it is used as a pointer and its own proof is not part of this item.\par\medskip
\ifdefined\hypertarget\hypertarget{sourcepointer-5}{}\fi
\noindent\textit{Source reference.} Original lemma 0788: sources/stacks-chow/tags.txt, source lines 9228--9228. Cross-chapter standard prerequisite cited by the author inside the counted span; it is used as a pointer and its own proof is not part of this item.\par\medskip
\ifdefined\hypertarget\hypertarget{sourcepointer-6}{}\fi
\noindent\textit{Source reference.} Original lemma 07B1: sources/stacks-chow/tags.txt, source lines 9326--9326. Cross-chapter standard prerequisite cited by the author inside the counted span; it is used as a pointer and its own proof is not part of this item.\par\medskip
\ifdefined\hypertarget\hypertarget{sourcepointer-7}{}\fi
\noindent\textit{Source reference.} Original proposition 0771: sources/stacks-chow/tags.txt, source lines 9186--9186. Cross-chapter standard prerequisite cited by the author inside the counted span; it is used as a pointer and its own proof is not part of this item.\par\medskip
\ifdefined\hypertarget\hypertarget{sourcepointer-8}{}\fi
\noindent\textit{Source reference.} Original lemma 0GR0: sources/stacks-chow/tags.txt, source lines 20880--20880. Cross-chapter standard prerequisite cited by the author inside the counted span; it is used as a pointer and its own proof is not part of this item.\par\medskip
\ifdefined\hypertarget\hypertarget{sourcepointer-9}{}\fi
\noindent\textit{Source reference.} Original reference 03DX: sources/stacks-chow/tags.txt, source lines 4436--4436. Cross-chapter standard prerequisite cited by the author inside the counted span; it is used as a pointer and its own proof is not part of this item.\par\medskip
\ifdefined\hypertarget\hypertarget{sourcepointer-10}{}\fi
\noindent\textit{Source reference.} Original reference 03DW: sources/stacks-chow/tags.txt, source lines 4435--4435. Cross-chapter standard prerequisite cited by the author inside the counted span; it is used as a pointer and its own proof is not part of this item.\par\medskip
\ifdefined\hypertarget\hypertarget{sourcepointer-11}{}\fi
\noindent\textit{Source reference.} Original reference 01XB: sources/stacks-chow/tags.txt, source lines 2512--2512. Cross-chapter standard prerequisite cited by the author inside the counted span; it is used as a pointer and its own proof is not part of this item.\par\medskip
\begin{lemma}
\label{lemma-quasi-coherent-weak-serre}
Let $\mathcal{X}$ be an algebraic stack.
\begin{enumerate}
\item $\QCoh(\mathcal{O}_{\mathcal{X}_{lisse,\etale}})$
is a weak Serre subcategory of
$\textit{Mod}(\mathcal{O}_{\mathcal{X}_{lisse,\etale}})$.
\item $\QCoh(\mathcal{O}_{\mathcal{X}_{flat,fppf}})$
is a weak Serre subcategory of
$\textit{Mod}(\mathcal{O}_{\mathcal{X}_{flat,fppf}})$.
\end{enumerate}
\end{lemma}

\begin{proof}
We will verify conditions (1), (2), (3), (4) of
Homology, Lemma \ref{homology-lemma-characterize-weak-serre-subcategory}.

\medskip\noindent
Since $0$ is a quasi-coherent module on any ringed site we see that (1)
holds.

\medskip\noindent
By definition $\QCoh(\mathcal{O})$
is a strictly full subcategory $\textit{Mod}(\mathcal{O})$, so (2) holds.

\medskip\noindent
Let $\varphi : \mathcal{G} \to \mathcal{F}$ be a morphism of quasi-coherent
modules on $\mathcal{X}_{lisse,\etale}$ or $\mathcal{X}_{flat,fppf}$.
We have $g^*g_!\mathcal{F} = \mathcal{F}$ and similarly for
$\mathcal{G}$ and $\varphi$, see Lemma \ref{lemma-lisse-etale-modules}.
By Lemma \ref{lemma-shriek-quasi-coherent}
we see that $g_!\mathcal{F}$ and $g_!\mathcal{G}$ are quasi-coherent
$\mathcal{O}_\mathcal{X}$-modules. By Sheaves on Stacks, Lemma
\ref{stacks-sheaves-lemma-quasi-coherent-algebraic-stack}
we have that $\Coker(g_!\varphi)$ is a quasi-coherent
module on $\mathcal{X}$ (and the cokernel in the category
of quasi-coherent modules on $\mathcal{X}$).
Since $g^*$ is exact (see Lemma \ref{lemma-lisse-etale})
$g^*\Coker(g_!\varphi) = \Coker(g^*g_!\varphi) = \Coker(\varphi)$
is quasi-coherent too (see Lemma \ref{lemma-quasi-coherent}).
By Proposition \ref{proposition-loc-qcoh-flat-base-change}
the kernel $\Ker(g_!\varphi)$ is in
$\textit{LQCoh}^{fbc}(\mathcal{O}_\mathcal{X})$.
Since $g^*$ is exact, we have
$g^*\Ker(g_!\varphi) = \Ker(g^*g_!\varphi) = \Ker(\varphi)$.
Since $g^*$ maps objects of
$\textit{LQCoh}^{fbc}(\mathcal{O}_\mathcal{X})$ to quasi-coherent
modules by Lemma \ref{lemma-quasi-coherent} we conclude that
$\Ker(\varphi)$ is quasi-coherent as well. This proves (3).

\medskip\noindent
Finally, suppose that
$$
0 \to \mathcal{F} \to \mathcal{E} \to \mathcal{G} \to 0
$$
is an extension of $\mathcal{O}_{\mathcal{X}_{lisse,\etale}}$-modules
(resp.\ $\mathcal{O}_{\mathcal{X}_{flat,fppf}}$-modules) with $\mathcal{F}$
and $\mathcal{G}$ quasi-coherent. To prove (4) and finish the proof
we have to show that
$\mathcal{E}$ is quasi-coherent on $\mathcal{X}_{lisse,\etale}$
(resp.\ $\mathcal{X}_{flat,fppf}$). Let $U$
be an object of $\mathcal{X}_{lisse,\etale}$
(resp.\ $\mathcal{X}_{flat,fppf}$; we think of $U$ as a scheme smooth
(resp.\ flat) over $\mathcal{X}$. We have to show that the restriction
of $\mathcal{E}$ to $U_{lisse,\etale}$ (resp.\ $=U_{flat,fppf}$)
is quasi-coherent. Thus we may assume
that $\mathcal{X} = U$ is a scheme. Because $\mathcal{G}$
is quasi-coherent on $U_{lisse,\etale}$ (resp.\ $U_{flat,fppf}$),
we may assume, after replacing $U$ by the members
of an \'etale (resp.\ fppf) covering, that $\mathcal{G}$ has a
presentation
$$
\bigoplus\nolimits_{j \in J} \mathcal{O} \longrightarrow
\bigoplus\nolimits_{i \in I} \mathcal{O} \longrightarrow
\mathcal{G} \longrightarrow 0
$$
on $U_{lisse,\etale}$ (resp.\ $U_{flat,fppf}$) where $\mathcal{O}$
is the structure sheaf on the site. We may also assume $U$ is affine.
Since $\mathcal{F}$ is quasi-coherent, we have
$$
H^1(U_{lisse,\etale}, \mathcal{F}) = 0,
\quad\text{resp.}\quad
H^1(U_{flat,fppf}, \mathcal{F}) = 0
$$
Namely, $\mathcal{F}$ is the pullback of a quasi-coherent module
$\mathcal{F}'$ on the big site of $U$
(by Lemma \ref{lemma-quasi-coherent}), cohomology
of $\mathcal{F}$ and $\mathcal{F}'$ agree (by
Lemma \ref{lemma-lisse-etale-cohomology}), and we know that
the cohomology of $\mathcal{F}'$ on the big site of the affine scheme $U$
is zero (to get this in the current situation you have to
combine Descent, Propositions
\ref{descent-proposition-equivalence-quasi-coherent} and
\ref{descent-proposition-same-cohomology-quasi-coherent} with
Cohomology of Schemes, Lemma
\ref{coherent-lemma-quasi-coherent-affine-cohomology-zero}).
Thus we can lift the map
$\bigoplus_{i \in I} \mathcal{O} \to \mathcal{G}$
to $\mathcal{E}$. A diagram chase shows that we obtain
an exact sequence
$$
\bigoplus\nolimits_{j \in J} \mathcal{O} \to
\mathcal{F} \oplus \bigoplus\nolimits_{i \in I} \mathcal{O}
\to
\mathcal{E} \to 0
$$
By (3) proved above, we conclude that $\mathcal{E}$ is quasi-coherent
as desired.
\end{proof}

\end{document}
